Battery-powered legged robots draw drive currents through one shared direct-current bus whose voltage falls under load, yet locomotion policies are trained assuming an ideal supply per joint. As the engineering application, we model one leg of a Unitree Go1-size quadruped, with two drive models checked against a switching-level simulation and a battery fitted to public cell data. As the artificial-intelligence contribution, a barrier-based filter between a reinforcement-learning policy and the drives reduces discrete-time barrier conditions on bus voltage and battery current to one power budget per step and projects the currents onto it. We train soft actor-critic policies with the filter in the loop and compare them with baseline policies (no filter, no penalty), test-time filtering, reward penalties and a Lagrangian method. Without the filter, every policy that hopped breached a limit; at the lowest state of charge, baseline policies spent \res{base_uv_pct_soc0.2}\% of each episode below the voltage limit. With exact battery parameters, every policy run with the filter met both limits. Policies trained with it never fell in any of their \res{filt_neps_all} evaluation episodes, including all stress tests. Policies filtered only at test time sometimes fell: baseline policies in every stress-test setting, and one of three penalty-trained seeds in three settings, though penalty policies hopped higher. A battery resistance below the model's caused current-limit breaches, which an extended Kalman filter estimating the resistance removed; measurement noise far above its assumed level still caused breaches. The guarantee assumes known battery electromotive force and resistance.
